% Part: incompleteness
% Chapter: representability-in-q
% Section: undecidability

\documentclass[../../../include/open-logic-section]{subfiles}

\begin{document}

\olfileid{inc}{req}{und}
\olsection{अनिर्णेयता}

एखाद्या उपपत्तीला $\Th{T}$ \emph{अनिर्णेय} म्हणतो, जर $\Th{T}$ च्या
भाषेतील कोणतेही वाक्य~$!A$ दिले असता ``$\Th{T}$ हे~$!A$
सिद्ध करते का?'' या प्रश्नाचे योग्य उत्तर सांत पायऱ्यांत आणि
नेहमी देणारी संगणकीय पद्धत अस्तित्वात नसेल. म्हणून अंकगणिताच्या
भाषेतील वाक्य~$!A$ दिल्यावर $\Th{Q} \Proves !A$ आहे की नाही,
हे ठरवणारी संगणकीय पद्धत असेल तेव्हाच आणि तरच $\Th{Q}$
निर्णेय असेल. हे अधिक नेमके करण्यासाठी विचारू: $\Prov[\Th{Q}](y)$
हा संबंध पुनरावर्ती आहे का? येथे तो~$y$ साठी तेव्हाच सत्य
असतो, जेव्हा $y$ ही~$\Th{Q}$ मध्ये सिद्ध होणाऱ्या वाक्याची
गोडेल संख्या असते. उत्तर नाही.

\begin{thm}
$\Th{Q}$ अनिर्णेय आहे; म्हणजे पुढील संबंध
\[
\Prov[\Th{Q}](y) \defiff \fn{Sent}(y) \land
\lexists[x][\Prf[\Th{Q}](x, y)]
\]
पुनरावर्ती नाही.
\end{thm}

\begin{proof}
तो पुनरावर्ती आहे असे समजा. मग थांबण्याची समस्या अशी
सोडवता येईल: $e$ आणि $n$ दिले असता,
$\cfind{e}(n) \fdefined$ तेव्हाच आणि तरच असा एखादा~$s$
असतो की $T(e, n, s)$, हे आपल्याला माहीत आहे. येथे $T$ हे
\olref[cmp][rec][nft]{thm:kleene-nf} मधील क्लिनीचे विधेय
आहे. $T$ आदिम पुनरावर्ती असल्याने त्याचे~$\Th{Q}$ मध्ये
$!B_T$ या सूत्राने निरूपण होते; म्हणजे
$\Th{Q}
\Proves !B_T(\num{e}, \num{n}, \num{s})$ तेव्हाच आणि तरच
$T(e, n, s)$. जर $\Th{Q}
\Proves !B_T(\num{e}, \num{n}, \num{s})$ असेल, तर
$ \Th{Q} \Proves
\lexists[y][!B_T(\num{e}, \num{n}, y)]$ हेसुद्धा मिळते.
असा एकही $s$ नसेल, तर प्रत्येक~$s$ साठी
$\Th{Q} \Proves \lnot !B_T(\num{e}, \num{n}, \num{s})$.
पण $\Th{Q}$ ही $\omega$-सुसंगत आहे: प्रत्येक~$n \in \Nat$
साठी $\Th{Q}
\Proves \lnot !A(\num{n})$ असेल, तर
$\Th{Q}
\Proves/ \lexists[y][!A(y)]$. कारण $\Th{Q}$ ची स्वयंसिद्धके
मानक प्रतिमानात~$\Struct{N}$ खरी आहेत. म्हणून
$\Th{Q}
\Proves/ \lexists[y][!B_T(\num{e}, \num{n}, y)]$.
दुसऱ्या शब्दांत, $\Th{Q}
\Proves \lexists[y][!B_T(\num{e}, \num{n}, y)]$ तेव्हाच आणि
तरच असा एखादा $s$ असतो की $T(e, n, s)$; म्हणजे तेव्हाच
आणि तरच $\cfind{e}(n) \fdefined$. $e$ आणि~$n$ पासून
$\Gn{\lexists[y][!B_T(\num{e},
    \num{n}, y)]}$ आपण संगणित करू शकतो; तसे करणारे
$g(e, n)$ हे आदिम पुनरावर्ती फलन असू द्या. मग
\[
h(e, n) =
\begin{cases}
1 & \text{जर $\Prov[\Th{Q}](g(e, n))$ असेल}\\
0 & \text{अन्यथा}.
\end{cases}
\]
$\Prov[\Th{Q}]$ पुनरावर्ती असेल, तर यावरून $h$ सुद्धा
पुनरावर्ती ठरेल. पण~$h$ हे
\olref[cmp][rec][hlt]{thm:halting-problem} नुसार पुनरावर्ती नाही.
म्हणून $\Prov[\Th{Q}]$ हाही पुनरावर्ती नाही.
\end{proof}

\begin{cor}
प्रथम-क्रम तर्कशास्त्र अनिर्णेय आहे.
\end{cor}

\begin{proof}
प्रथम-क्रम तर्कशास्त्र निर्णेय असते, तर~$\Th{Q}$ मधील
सिद्धतायोग्यताही निर्णेय असती: कारण $\Th{Q} \Proves !A$
तेव्हाच आणि तरच $\Proves !T \lif !A$; येथे $!T$ हा
$\Th{Q}$ च्या सांत स्वयंसिद्धकांचा संयोग आहे.
\end{proof}

\end{document}
